%%%%

%%%   Самарский государственный технический университет

%%%   Кафедра прикладной математики и информатики

%%%

%%%   Преамбула для статьи в журнал "Вестник Самарского государственного

%%%   технического университета. Серия: «Физико-математические науки»"

%%%   Ориентирована под MikTEx 2.4 for Win32

%%%

%%%   Хотя сам журнал собирается под GNU/Linux на дистрибутиве TeXLive



\documentclass[11pt,a4paper,twoside]{article}



\usepackage[cp1251]{inputenc}

\usepackage[T2A]{fontenc}

\usepackage[english,russian]{babel}

\usepackage{samgtubib}

\usepackage[dvips]{graphicx}

\usepackage[multiple]{footmisc}



\usepackage{samgtu}

\renewcommand{\VestnikYear}{2009} % Год издания Вестника

\renewcommand{\VestnikNumber}{2\,(19)} % Номер Вестника

\renewcommand{\VestnikMonth}{Ноябрь} % Месяц выхода Вестника

\renewcommand{\VestnikData}{01.11.09} % Дата подписания Вестника в печать

\renewcommand{\VestnikUPL}{26,64} % Усл. печ. л.

\renewcommand{\VestnikUIL}{25,73} % Уч.-изд. л.

\renewcommand{\VestnikRegNumber}{479/09} % Регистрационный номер

\renewcommand{\VestnikOrder}{994} % Номер заказа







%

%\usepackage{pst-barcode}

%\usepackage{auto-pst-pdf}

%%

%%\newcommand{\totop}[1]{\raisebox{6pt}[1pt][2pt]{#1}} 

%%\newcommand{\tottop}[1]{\raisebox{12pt}[1pt][2pt]{#1}} 

%%\newcommand{\totttop}[1]{\raisebox{18pt}[1pt][2pt]{#1}} 

%%\newcommand{\tottttop}[1]{\raisebox{24pt}[1pt][2pt]{#1}} 

%%\newcommand{\totttttop}[1]{\raisebox{30pt}[1pt][2pt]{#1}} 

%%\newcommand{\tobot}[1]{\raisebox{-6pt}[1pt][2pt]{#1}} 

%%\newcommand{\tobbot}[1]{\raisebox{-12pt}[1pt][2pt]{#1}} 

%%\newcommand{\tobbbot}[1]{\raisebox{-18pt}[1pt][2pt]{#1}}

%

%\graphicspath{{./00PIC/}}

%

%%\samgtupubl % для генерации pdf, отдаваемого в типографию СамГТУ

%\pdfcrop % обрезка pdf для размещения на math-net.ru

%\draft % На корректуре

%\filllastpage % Последняя страница не заполнена не полностью

%%\briefcomm % Статья является кратким сообщением



\vsgtu{1438}



\FirstPage{149}

\LastPage{157}

%

%

%\begin{document}



\hfill \qrcode[hyperlink,height=.8in]{http://dx.doi.org/10.14498/vsgtu1438}

\vspace{-.8in}





%

%

%\hfill \begin{pspicture}(1in,1in)

%\psbarcode{http://dx.doi.org/10.14498/1438}{}{qrcode}

%\end{pspicture}

%\vspace{-1in}





\udc{512.14, 519.117}

\msc{05A10, 05A19}



\rutitle{Суммирование на основе комбинаторного представления одинаковых степеней}

{Суммирование на основе комбинаторного представления\dots}



\ruauthor{А.~И.~Никонов}

{Н\,и\,к\,о\,н\,о\,в~А.~И.}





\ruaffil{\rusamgtu.}







\ruauthordetails{1}{\fullauthor{\rupid{39206}{Александр Иванович Никонов}} \regalia{д.т.н., проф.; \mem{nikonovai@mail.ru}}\position{профессор} \dept{каф. электронных систем и информационной безопасности}}





\rukeywords{сумма одинаковых степеней, степень, геометрическая прогрессия, арифметико"=геометрическая прогрессия, комбинированная прогрессия, матрица}



\ruabstract{В статье рассматривается вывод комбинаторных выражений для сумм членов нескольких последовательностей. Вывод производится на основе комбинаторного представления суммы взвешенных одинаковых степеней. Суммированию подлежат взвешенные члены геометрической прогрессии, простых арифметико"=геометрической и комбинированной прогрессий. Одно из главных мест в данном выводе занимает представление членов каждой из указанных прогрессий в виде элементов матрицы. Строка этой матрицы формируется с использованием набора одинаковых степеней с заданным весовым коэффициентом. Кроме того, в работе представлены формулы комбинаторных тождеств с участием свободных компонентов сумм одинаковых степеней, а также отдельной степени "--- члена последовательности одинаковых степеней или геометрической прогрессии. Все представленные формулы имеют общую основу "--- компоненты сумм одинаковых степеней.}



\entitle{Summation on the basis of combinatorial representation of equal powers}



\enauthor{A.~I.~Nikonov}{N\,i\,k\,o\,n\,o\,v~A.~I.}



\enaffil{\ensamgtu.}



\enauthordetails{1}{\fullauthor{\enpid{39206}{Alexander I. Nikonov}}

\regalia{Dr. Tech. Sci.; \mem{nikonovai@mail.ru}}\position{Professor} \dept{Dept. of

Electronic Systems and Information Security}}



\enabstract{In the paper the conclusion of combinatorial expressions for the sums of members of several sequences is considered. Conclusion is made on the basis of combinatorial representation of the sum of the weighted equal powers. The weighted members of a geometrical progression, the simple arithmetic-geometrical and combined progressions are subject to summation. One of principal places in the given conclusion occupies representation of members of each of the specified progressions in the form of matrix elements. The row of this matrix is formed with use of a gang of equal powers with the set weight factor. Besides, in work formulas of combinatorial identities with participation of free components of the sums of equal powers, and also separate power-member of sequence of equal powers or a geometrical progression are presented. All presented formulas have the general basis-components of the sums of equal powers. }



\enkeywords{sum of equal powers, power, geometrical

progression, simple arithmetic-geometrical progression, 

simple combined progression, matrix}



\date{27/V/2015}

\revisiondate{23/VII/2015}

\accepteddate{08/VIII/2015}



\makerutitle





Комбинаторным представлением какого"=либо математического соотношения  мы называем такое его представление, которое производится с применением объектов комбинаторики, например, сочетаний и выражений на их основе~[\citen{nikonov:bib1}--\citen{nikonov:bib4}]. В комбинаторном представлении математических выражений можно выделить представление сумм одинаковых степеней~\cite{nikonov:bib5,nikonov:bib6}, в том числе сумм взвешенных одинаковых степеней~\cite{nikonov:bib7,nikonov:bib8}. Явное комбинаторное представление суммы взвешенных одинаковых степеней (СВОС) выглядит следующим образом~\cite{nikonov:bib8}:

\begin{equation}

\label{nikonov:eq1}

\Phi_{\rm oc}(p,\nu)=\sum\limits_{l=1}^p b_l l^\nu=\sum \limits_{\iota=1}^{\max \iota} \left(\sum \limits_{q=1}^\iota (-1)^{\iota+q} C^{q}_\iota q^\nu \right) \left(\sum_{k=0}^{p-\iota} C^\iota_{p-k}b_{p-k} \right),

\end{equation}

где $p$ "--- число одночленов, в данном случае с весовыми коэффициентами вида $b_l \in \mathbb{R}$, $p \in \mathbb{N}$; $\nu$ "--- одинаковая степень натуральных чисел вида $l$, $\nu \in \mathbb{N}$; $\max \iota=\min (p,\nu)$.



Целью настоящей статьи является получение комбинаторных выражений для сумм членов нескольких последовательностей взвешенного типа, которое производится на основе комбинатроного представления СВОС вида \eqref{nikonov:eq1}. Этими последовательностями, рассматриваемыми в первой части данной статьи, являются геометрическая прогрессия (ГП), простая арифметико"=геометрическая прогрессия (АГП), простая комбинированная прогрессия (КП)~[\citen{nikonov:bib9}--\citen{nikonov:bib11}]. Они играют, в частности, заметную роль в описании процессов шифрования с участием сумм со слагаемыми"=произведениями свободных и весовых компонентов~\cite{nikonov:bib11}. Во второй части нашей статьи рассматривается вывод двух комбинаторных тождеств с участием свободных компонентов СВОС.



Целочисленные аргументы, рассматриваемые в настоящем исследовании, позволяют обеспечить как приближённую оценку суммируемых членов прогрессий, так и их точную оценку в случае наличия рациональных значений знаменателей данных сумм. Исследование второго из указанных выше комбинаторных тождеств (вторая часть нашей статьи) избавлено от целочисленного характера.



В дальнейшем для краткости мы будем иногда именовать каждую из рассматриваемых прогрессий взвешенной. Сумма членов любой из данных прогрессий относится к многочленам одной переменной, но имеет и специфическое вышеприведённое название, основанное на названии соответствующей последовательности~\cite{nikonov:bib11}.



Сумма взвешенной геометрической прогрессии с целочисленным знаменателем $\rho$ "--- числом одночленов в данном случае с весовыми коэффициентами вида $b_{\nu l}=b_\nu$ выглядит следующим образом:

$$

\Phi_{\text{гп}}(\rho,\max \nu)=\sum \limits_{\nu=\min \nu}^{\max \nu} b_\nu \rho^\nu, \quad b_\nu \in \mathbb{R}.

$$



Для определённости примем $\min \nu=1$; $\max \nu=\tilde{p} \in \mathbb{N}$.



Введём в рассмотрение матрицу взвешенных членов, см. рис.~\ref{nik:fig:1}. Её строка состоит из взвешенных членов последовательности

$$

\left(b_\nu l^\nu, \quad l=1,\ldots, p \right).

$$



\begin{figure}[h!]\centering\footnotesize

\includegraphics{nikonov1}



\caption{Матрица взвешенных членов для ГП \label{nik:fig:1}}

\smallskip



[Figure~\ref{nik:fig:1}. A suspended members matrix for geometric progression]



\end{figure}



Суммы членов взвешенной геометрической прогрессии для знаменателей $p$, $p-1$ и степенного показателя

$$

\nu=\tilde{p}-k, \quad k=0, 1, \ldots,\tilde{p}-1

$$

будут соответствовать набору элементов СВОС произвольной строки рассматриваемой матрицы, а также набору указанных элементов, простирающихся лишь до двойной линии, проведённой в этой матрице.



Аналитически данное рассуждение может быть выражено таким образом:

\begin{equation}

\label{nikonov:eq2}

\Phi_{\text{гп}}(p,\tilde{p})=\sum \limits_{\nu=1}^{\tilde{p}} b_\nu p^\nu, \quad \Phi_{\text{гп}} (p-1,\tilde{p}) = \sum \limits_{\nu=1}^{\tilde{p}-1} b_\nu (p-1)^\nu,

\end{equation}

причем

\begin{equation}

\label{nikonov:eq3}

\Delta \Phi_{\text{гп}}(p,\tilde{p})=\Phi_{\rm oc} (p,\nu) -\Phi_{\rm oc}(p-1,\nu)=b_\nu p^\nu.

\end{equation}



Из формулы \eqref{nikonov:eq3} находим:

\begin{equation}

\label{nikonov:eq4}

\Phi_{\text{гп}}(p,\tilde{p})=\sum \limits_{\nu=1}^{\tilde{p}} b_\nu p^\nu= \sum \limits_{\nu=1}^{\tilde{p}} \Delta \Phi_{\text{гп}} (p,\tilde{p}).

\end{equation}



Здесь величины $\Phi_{\rm oc}(p,\nu)$, $\Phi_{\rm oc}(p-1,\nu)$ "--- это суммы взвешенных одинаковых степеней с числом элементов соответственно $p$, $p-1$ и одинаковым показателем $\nu$.



Подставив в выражение \eqref{nikonov:eq4} формулу \eqref{nikonov:eq1} с параметрами $p$, $p-1$, получим следующее:

\begin{equation}

\label{nikonov:eq5}

\begin{aligned}

&\Phi_{\text{гп}} (p,\tilde{p})=\sum\limits_{\nu=1}^{\tilde{p}} b_\nu A_\nu;\\

&A_\nu=\sum \limits_{\iota=1}^{\max \iota} \left(\sum \limits_{q=1}^\iota (-1)^{\iota +q} C^q_\iota q^\nu \right) C^\iota_p.

\end{aligned}

\end{equation}



Совокупность выражений \eqref{nikonov:eq5} и есть искомая комбинаторная формула суммы взвешенных членов ГП, соответствующая представлению \eqref{nikonov:eq2}.



Использовав теперь непосредственную подстановку в формулу \eqref{nikonov:eq1} значения $l=0$, будем иметь

$$

\Delta \Phi_{\text{гп}} (p, \nu=0)=\Delta \Phi_{\text{гп}} (p, 0)\big|_{\iota=0,q=0}=b_0 (-1)^0 C_0^0\delta_{0\nu}^0 C_p^0=b_0,

$$

где $b_0$ "--- задаваемый весовой коэффициент с нулевым индексом; $\delta_{0\nu}$ "--- символ Кронекера, вводимый сюда, поскольку член $0^0$ принято считать не имеющим смысла.



Тогда формула суммы взвешенных членов ГП \eqref{nikonov:eq5} получает модифицированный вид

\begin{equation}

\label{nikonov:eq6}

\begin{aligned}

&\Phi_{\text{гп}} (p,\tilde{p})=\sum\limits_{\nu=0}^{\tilde{p}} b_\nu A_{\nu \delta};\\

&A_{\nu\delta}=\sum \limits_{\iota=0}^{\max \iota} \left(\sum \limits_{q=0}^\iota (-1)^{\iota +q} C^q_\iota(q+\delta_{0\nu})^\nu \right) C^\iota_p.

\end{aligned}

\end{equation}



Этот вид соответствует такому представлению ГП, как

$$

\Phi_{\text{гп}} (p,\tilde{p})=\sum\limits_{\nu=0}^{\tilde{p}} b_\nu p^\nu,

$$

что, в свою очередь, связано с добавлением в матрицу взвешенных членов на рис.~\ref{nik:fig:1} новой нижней строки

$$

\left(b_0 l^0=b_0, \ldots l=1,\ldots,p \right).

$$



Как видим, выражение \eqref{nikonov:eq6} позволяет расширить область действия формулы \eqref{nikonov:eq5}. Тем не менее, далее в настоящей работе случаи нулевого нижнего предела определяемых ниже сумм рассматриваться не будут, поскольку данные суммы имеют в таких случаях, в отличие от суммы взвешенной ГП, лишь нулевые значения.



Простая форма представления суммы членов взвешенной АГП имеет вид

\begin{equation}

\label{nikonov:eq7}

\Phi_{\text{агп}} (p,\tilde{p})=\sum\limits_{\nu=1}^{\tilde{p}} b_\nu \nu p^\nu=\sum \limits_{\nu=1}^{\tilde{p}} b_\nu' p^{\nu}; \quad b'_\nu =b_\nu \nu.

\end{equation}



Тогда искомая формула комбинаторного представления суммы

$\Phi_{\text{агп}}$, записанной в виде \eqref{nikonov:eq7}, будет

выражаться в соответствии с соотношением \eqref{nikonov:eq4} как

\begin{equation}

\label{nikonov:eq8}

\Phi_{\text{агп}}(p,\tilde{p})=\sum \limits_{\nu=1}^{\tilde{p}} b'_\nu A_\nu.

\end{equation}



Структуры выражений \eqref{nikonov:eq7} и \eqref{nikonov:eq2}

соответствуют друг другу, и ясно, что матрица взвешенных членов для

простой взвешенной АГП будет аналогична матрице для взвешенной ГП

с учетом, естественно, равенства \eqref{nikonov:eq8}. Этот учет

дает возможность применить к выражению \eqref{nikonov:eq8} формулу

СВОС для единичной одинаковой степени со значениями соответствующего коэффициента $\nu^1=1$ и с весовыми коэффициентами в $b_\nu A_\nu=b_\nu A_{\tilde{p}-k}$ при $\alpha_0^1(1)=1$:

\begin{gather}

\label{nikonov:eq9}

\Phi_{\text{агп}}=\sum \limits_{k=0}^{\tilde{p}-1}

C^1_{\tilde{p}-k} A_\nu, \quad \tilde{p}-k=\nu;\\

\label{nikonov:eq10}

\max \iota= \max \iota(p,\nu)=\min(p,\nu).

\end{gather}



Соотношения \eqref{nikonov:eq9}, \eqref{nikonov:eq10} дают искомое

комбинаторное представление суммы взвешенных членов простой АГП.



Перейдем теперь к исследованию выражения суммы взвешенных

членов простой КП:

\begin{equation}

\label{nikonov:eq11} \Phi_{\text{кп}} (p,p)=\sum

\limits_{l=1}^{\tilde{p}} b_l (dl)^{dl}, \quad  d \in \mathbb{N}.

\end{equation}



Рассмотрим рис.~\ref{nik:fig:2}, где изображено положение чисел $b_l (dl)^{dl}$

"--- компонентов выражения \eqref{nikonov:eq11}; слева направо

идет нарастание значений $l$:

$$

(dl)^{dl}=d^{dl} l^{dl}, \quad l=1,\ldots, \nu_d=\max l, \quad

\nu_d=1,\ldots,\max \nu_d=p.

$$





\begin{figure}[h!]\centering\footnotesize

\includegraphics{nikonov2}

\caption{Матрица взвешенных членов для простой КП \label{nik:fig:2}}



\smallskip



[Figure~\ref{nik:fig:2}. A suspended members matrix for easy combination progression]



\end{figure}



Комбинаторная формула суммы взвешенных членов простой КП выводится

аналогично формуле суммы взвешенных членов ГП \eqref{nikonov:eq5}

со следующими особенностями:

\begin{itemize}

\item[а)] $\nu=d\nu_d$; 

\item[б)] $p=\tilde{p}=\max

\nu_d$; 

\item[в)] в выводимом выражении присутствует множитель $d^\nu$.



\end{itemize}

Тогда имеем

\begin{equation}

\label{nikonov:eq12}

\begin{aligned}

&\Phi_{\text{кп}}(p,p)=\sum \limits_{\nu_d=1}^{p}

b_{\nu_d} d^{d\nu_d}A_{\nu_d}, \quad b_{\nu_d}=b_{\max l};\\

&A_{\nu_d}= \sum \limits_{\iota=1}^{\max \iota} \left( \sum

\limits_{q=1}^\iota (-1)^{\iota+q} C^q_\iota q^\nu \right)

C^\iota_{\nu_d};\\

&\max \iota = \min(\nu_d, \nu_d)=\nu_d.

\end{aligned}

\end{equation}



Слагаемые из \eqref{nikonov:eq12} располагаются в строке матрицы

на рис.~\ref{nik:fig:2} в её первой колонке и правее двойной линии. Переходя к

преобразованию \eqref{nikonov:eq12}, видим, что мы можем

воспользоваться здесь еще раз, как и при преобразовании формулы

\eqref{nikonov:eq11}, выражением \eqref{nikonov:eq5}:

\begin{equation}

\label{nikonov:eq13}

\begin{array}{c}

\displaystyle \Phi_{\text{кп}} (p',\nu_d)=\sum \limits_{\nu_d=1}^p b_{\nu_d} \left(\sum \limits_{\iota_1=1}^{\max \iota_1}

\left(\sum \limits_{q=1}^{\iota_1} (-1)^{\iota_1+q}C^q_{\iota_1} q^{\nu_d}\right) C^{\iota_1}_{p'} \right) A_{\nu_d}, \\

 p'=d^d, \quad \max \iota_1=\min \left(p',\nu_d \right).

\end{array}

\end{equation}



Формула \eqref{nikonov:eq13} и есть искомое комбинаторное выражение суммы взвешенных членов простой КП.



При $d=1$ и, соответственно, при $p'=1$, $\nu=\nu_d$ имеет место важный частный случай выражения \eqref{nikonov:eq13}:

\begin{equation}

\label{nikonov:eq14}

\begin{array}{c}

\displaystyle

\Phi_{\text{кп}} (p',\nu_d)= \Phi_{\text{кп}} (1,\nu_d)

=\sum \limits_{\nu_d=1}^p b_{\nu_d} \left(\sum \limits_{\iota_1=1}^{\max \iota_1}

\left(\sum \limits_{q=1}^{\iota_1} (-1)^{\iota_1+q}C^q_{\iota_1} q^{\nu_d}\right)C^{\iota_1}_{p'}  \right) A_{\nu_d}, \\

\max \iota_1=\min \left(p',\nu_d \right)=1.

\end{array}

\end{equation}



То есть исходя из соотношения \eqref{nikonov:eq14}, получаем

$$

\Phi_{\text{кп}}(1,\nu_d)=\sum \limits_{\nu_d=1}^p b_{\nu_d} A_{\nu_d}.

$$



В заключительной части нашей работы выведем сперва комбинаторное тождество, соответствующее матрице рис.~\ref{nik:fig:1}. Согласно выражению, аналогичному \eqref{nikonov:eq3},

при $b_\nu=1$ имеем

\begin{equation}

\label{nikonov:eq15}

\Delta \Phi_{\rm oc} (p,\nu)=p^\nu,

\end{equation}

но в то же время

\begin{equation}

\label{nikonov:eq16}

\Delta \Phi_{\rm oc} (p,\nu)= \sum \limits_{\iota=1}^{\max \iota} \alpha_0^\iota (\nu)  C^\iota_p, \quad \max \iota=\min (p,\nu),

\end{equation}

где $\alpha_0^\iota (\nu)$ "--- свободный компонент СВОС, причем с учетом формулы \eqref{nikonov:eq1}

$$

\alpha_0^\iota (\nu) = \sum \limits_{q=1}^\iota (-1)^{\iota+q} C^q_\iota q^\nu.

$$



Если принять $p\leq \nu$, то, учитывая \eqref{nikonov:eq15}, \eqref{nikonov:eq16}, получим следующее тождество для свободного компонента СВОС:

\begin{equation}

\label{nikonov:eq17}

\alpha_0^p (\nu)=p^\nu-\sum \limits_{\iota=1}^{\max \iota -1} \alpha_0^\iota (\nu) C^\iota_p.

\end{equation}



Используемое для подсчета свободных компонентов вида $\alpha_0^p

(\nu)$ тождество \eqref{nikonov:eq17} позволяет достаточно

экономно тратить вычислительные ресурсы на выполнение этого

подсчета. При $p\geq \nu$ число подсчитываемых свободных

компонентов не выходит (поскольку здесь $\max \iota=\nu$) за рамки

числа $\nu$, так что и в данном случае выражение

\eqref{nikonov:eq17} все равно остается справедливым.



Наконец, выявим соотношение, связывающее разности последовательных

сумм одинаковых степеней и членов ГП:

\begin{gather*}

\Delta \Phi_{\rm oc} (rl,l)= \Phi_{\rm oc} (rl,l)- \Phi_{\rm oc}

(r(l-1),l)=(rl)^l=r^ll^l;\\

\Delta \Phi_{\text{гп}} (rl,l)= \Phi_{\text{гп}} (rl,l)-

\Phi_{\text{гп}} (rl,l-1)=(rl)^l;\\

l \in \mathbb{N}, \quad r \in \mathbb{R}_+ \setminus[0,1).

\end{gather*}



Мы смогли взять здесь величины $\Delta \Phi_{\rm oc}$, $\Delta

\Phi_{\text{гп}}$, поскольку они именно являются одинаковыми

разностями $(rl)^l$ последовательных сумм одинаковых степеней и

членов ГП, см. рис.~\ref{nik:fig:1}. Величина $\Delta \Phi_{\rm oc}$ формируется

путем использования выражения \eqref{nikonov:eq1}, учета

количества одночленов $l$ вида $l^l$ и единого множителя $r^l$

этих одночленов. В основе формирования величины $\Delta

\Phi_{\text{гп}}$ лежит использование выражения

\eqref{nikonov:eq5} и учет числа $(rl)$, находящегося в основании

степени $(rl)^l$, причем количество этих степеней также равно $l$.



В соответствии с формулой \eqref{nikonov:eq17} имеем

\begin{gather*}

\Delta \Phi_{\rm oc} (rl,l)= (r^l) \sum \limits_{\iota=1}^{\max \iota}

\alpha _0^\iota (l) C_l^{\iota};\\

\Delta \Phi_{\text{гп}} (rl,l)= \sum \limits_{\iota=1}^{\max

\iota} \alpha _0^\iota (l) C_{rl}^{\iota};

\end{gather*}

$\max \iota$ в данном случае составляет $\min (rl, l)=l$;

\begin{equation}

\label{nikonov:eq18} \Delta \Phi_{\rm oc} (rl,l)=\Delta

\Phi_{\text{гп}} (rl,l)=(rl)^l.

\end{equation}



Применяя соотношение \eqref{nikonov:eq18}, получаем

\begin{equation}

\label{nikonov:eq19}

r^l \sum \limits_{\iota =1} ^{\max \iota} \alpha_0^{\iota} (l) C_{l}^{\iota}=\sum \limits_{\iota=1}^{\max \iota} \alpha^{\iota}_0 (l) C_{rl}^{\iota}.

\end{equation}



Рассмотрение структуры $C_{rl}^\iota$, содержащей целые числа вида

$\iota$, позволяет утверждать, что все используемые в найденном

выражении \eqref{nikonov:eq19} формулы являются конечными. Более

того, сомнений в истинности тождества \eqref{nikonov:eq19} для $r$

"--- целого положительного начинающегося с единицы числа,

естественно, не возникает; таких натуральных чисел $r$ "---

бесконечное количество. Но тогда, согласно принципу полиномиальной

аргументации \cite{nikonov:bib12}, данное выражение сохраняет свою истинность и для

любого вещественного $r$, лежащего в указанных выше границах.









\thanks{{\bf ORCID}



{\bf Александр Иванович Никонов:} \url{http://orcid.org/0000-0002-0943-6068}



}





\RusRef





\begin{thebibliography}{99}



\Bibitem{nikonov:bib1} 

\by Coolidge~J.~L. 

\paper  The story of the Binomial Theorem 

\jour Amer. Math. Monthly 

\vol 56 

\issue 3 

\pages 147--157

\yr 1949

\crossref{10.2307/2305028}



\Bibitem{nikonov:bib2}

\by  Lin Cong

\preprint On Bernoulli Numbers and Its Properties

\arxiv \href{http://arxiv.org/abs/math/0408082}{math/0408082} [math.HO]

\yr 2004

\totalpages 7 



\Bibitem{nikonov:bib3}

\by Benjamin~A.~T., Plott~S.~S.

\paper A combinatorial approach to Fibonomial coefcients

\jour Fibonacci Quarterly

\yr 2008/2009

\vol 46/47

\issue 1

\pages 7--9

\elink URL: \url{http://www.fq.math.ca/Papers1/46_47-1/Benjamin_11-08.pdf} \Russian (дата обращения: 08.08.2015)



\Bibitem{nikonov:bib4}

\by Benjamin~A.~T., Greg~O.~P., Quinn~J.~J.

\paper A Stirling Encounter with Harmonic Numbers

\jour Mathematics Magazine

\vol 75

\issue 2

\yr 2002

\pages 95--103

\crossref{10.2307/3219141}



\Bibitem{nikonov:bib5} 

\by Edwards~A.~W.~F. 

\paper Sums of Powers of Integers: A Little of the History 

\jour The Mathematical

Gazette \vol 66 \issue 435 \yr 1982 \pages 22--28

\crossref{10.2307/3617302}



\Bibitem{nikonov:bib6}

\by Beery~J.

\preprint Sums of Powers of Positive Integers - Introduction

\preprintinfo Convergence (July 2010)

\yr 2010

\crossref{10.4169/loci003284}



\RBibitem{nikonov:bib7}

\by Никонов~А.~И.

\paper Преобразование суммы взвешенных степеней натуральных чисел с~одинаковыми показателями

\jour Вестн. Сам. гос. техн. ун-та. Сер. Физ.-мат. науки

\yr 2010

\issue 1(20)

\pages 258--262

\mathnet{http://mi.mathnet.ru/vsgtu751}

\crossref{10.14498/vsgtu751}



\RBibitem{nikonov:bib8}

\by Никонов~А.~И.

\paper Приведение суммы взвешенных одинаковых степеней к~явному комбинаторному представлению

\jour Вестн. Сам. гос. техн. ун-та. Сер. Физ.-мат. науки

\yr 2012

\issue 3(28)

\pages 163--169

\mathnet{http://mi.mathnet.ru/vsgtu1099}

\crossref{10.14498/vsgtu1099}



\Bibitem{nikonov:bib9} 

\elink URL: \url{http://www.trans4mind.com/personal_development/mathematics/series/airthmeticGeometricSeries.htm} \Russian (дата обращения: 08.08.2015)

\preprint  Arithmetic and Geometric Series in Ken Ward's Mathematics Pages







\Bibitem{nikonov:bib10}

\by Riley~K.~F., Hobson~M.~P., Bence~S.~Y.

\book Mathematical Methods for Physics and Engi\-nee\-ring

\publaddr Cambridge

\publ Cambridge University Press

\yr 2006, 

%\totalpages 

xxiv+1232~pp

\crossref{10.1017/CBO9781139164979}



\RBibitem{nikonov:bib11}

\by Никонов~А.~И.

\paper Повышение эффективности шифрования на основе суммирования произведений

\jour Вестн. Сам. гос. техн. ун-та. Сер. Физ.-мат. науки

\yr 2014

\issue 2(35)

\pages 199--207

\mathnet{http://mi.mathnet.ru/vsgtu1316}

\crossref{10.14498/vsgtu1316}



\Bibitem{nikonov:bib12}

\by Graham~R.~L., Knuth~D.~E., Patashnik~O.

\book Concrete mathematics. A foundation for computer science

\publ Addison-Wesley Publishing Company

\publaddr Reading, MA

\yr 1994

\totalpages xiv+657 





\end{thebibliography}



\RuDate



\bigskip



\makeentitle







\thanks{{\bf ORCID}



{\bf Alexander I. Nikonov:} \url{http://orcid.org/0000-0002-0943-6068}



}





\EngRef





\begin{thebibliography}{99}



\Bibitem{nikonov:bib1:eng} 

\by Coolidge~J.~L. 

\paper  The story of the Binomial Theorem 

\jour Amer. Math. Monthly 

\vol 56 

\issue 3 

\pages 147--157

\yr 1949

\crossref{10.2307/2305028}



\Bibitem{nikonov:bib2:eng}

\by  Lin Cong

\preprint On Bernoulli Numbers and Its Properties

\arxiv \href{http://arxiv.org/abs/math/0408082}{math/0408082} [math.HO]

\yr 2004

\totalpages 7 



\Bibitem{nikonov:bib3:3:eng}

\by Benjamin~A.~T., Plott~S.~S.

\paper A combinatorial approach to Fibonomial coefcients

\jour Fibonacci Quarterly

\yr 2008/2009

\vol 46/47

\issue 1

\pages 7--9

\elink  Retrieved from   \url{http://www.fq.math.ca/Papers1/46_47-1/Benjamin_11-08.pdf} (August 08, 2015)



\Bibitem{nikonov:bib4:eng}

\by Benjamin~A.~T., Greg~O.~P., Quinn~J.~J.

\paper A Stirling Encounter with Harmonic Numbers

\jour Mathematics Magazine

\vol 75

\issue 2

\yr 2002

\pages 95--103

\crossref{10.2307/3219141}



\Bibitem{nikonov:bib5:eng} 

\by Edwards~A.~W.~F. 

\paper Sums of Powers of Integers: A Little of the History 

\jour The Mathematical

Gazette \vol 66 \issue 435 \yr 1982 \pages 22--28

\crossref{10.2307/3617302}



\Bibitem{nikonov:bib6:eng}

\by Beery~J.

\preprint Sums of Powers of Positive Integers - Introduction

\preprintinfo Convergence (July 2010)

\yr 2010

\crossref{10.4169/loci003284}



\Bibitem{nikonov:bib7:eng}

\by Nikonov~A.~I.

\paper Converting the Sum of Weighted Degrees of Natural Numbers with the Same Parameters

\jour Vestn. Samar. Gos. Tekhn. Univ. Ser. Fiz.-Mat. Nauki  \rm [J. Samara State Tech. Univ., Ser.

Phys. \& Math. Sci.]

\yr 2010

\issue 1(20)

\pages 258--262

\crossref{10.14498/vsgtu751}



\Bibitem{nikonov:bib8:en}

\by Nikonov~A.~I.

\paper Reduction of the sum of the weight equal powers to explicit combinatorial representation

\jour Vestn. Samar. Gos. Tekhn. Univ. Ser. Fiz.-Mat. Nauki \rm [J. Samara State Tech. Univ., Ser.

Phys. \& Math. Sci.]

\yr 2012

\issue  3(28)

\pages 163--169

\crossref{10.14498/vsgtu1099}



\Bibitem{nikonov:bib9:eng} 

\elink  Retrieved from  \href{http://www.trans4mind.com/personal_development/mathematics/series/airthmeticGeometricSeries.htm}{\tt http://www.trans4mind.com/personal\_development/mathematics/series/airthmetic\-GeometricSeries.htm}

(August 08, 2015)

\preprint  Arithmetic and Geometric Series in Ken Ward's Mathematics Pages







\Bibitem{nikonov:bib10:eng}

\by Riley~K.~F., Hobson~M.~P., Bence~S.~Y.

\book Mathematical Methods for Physics and Engi\-nee\-ring

\publaddr Cambridge

\publ Cambridge University Press

\yr 2006, 

%\totalpages 

xxiv+1232~pp

\crossref{10.1017/CBO9781139164979}



\Bibitem{nikonov:bib11:eng}

\by Nikonov~A.~I.

\paper Rising of Efficiency of Enciphering on the Basis of~Summation of Products

\jour Vestn. Samar. Gos. Tekhn. Univ. Ser. Fiz.-Mat. Nauki  \rm [J. Samara State Tech. Univ., Ser.

Phys. \& Math. Sci.]

\yr 2014

\issue 2(35)

\pages 199--207

\crossref{10.14498/vsgtu1316}



\Bibitem{nikonov:bib12:eng}

\by Graham~R.~L., Knuth~D.~E., Patashnik~O.

\book Concrete mathematics. A foundation for computer science

\publ Addison-Wesley Publishing Company

\publaddr Reading, MA

\yr 1994

\totalpages xiv+657 





\end{thebibliography}





\EngDate



\newpage

\Russian





%\end{document}

